% Generated by TeXFlow — do not edit.
% Changes will be overwritten on next compile.
% Use the TeXFlow edit() and layout() tools to modify the document.
\documentclass[10pt, conference, letterpaper]{IEEEtran}

\IEEEoverridecommandlockouts
\markboth{Syeda Salsabil Islam Ariya}{Quantization and Sampling Effects in a Digitally Controlled PID Buck-Boost Converter: A Simulation Study}
\title{Quantization and Sampling Effects in a Digitally Controlled PID Buck-Boost Converter: A Simulation Study}
\author{Syeda Salsabil Islam Ariya and Jasimul Islam Chowdhury and Hussain Touhid Siddiquee}

\usepackage{amsmath}
\usepackage{amssymb}
\usepackage{booktabs}
\usepackage{graphicx}
\usepackage{url}

\begin{document}
\maketitle

\section{Introduction}

Digital control loops are now standard in switching power converters, driven by programmability, immunity to analog component drift, and integration with digital loads \cite{maksimovic2004,corradini2015}. The complete digital loop contains two unavoidable quantizers: the analog-to-digital converter (ADC) that samples the output voltage, and the digital pulse-width modulator (DPWM) that resolves the duty ratio. Their resolutions, together with the sampling period of the controller, jointly determine regulation quality. Classic work established that coarse quantization causes steady-state limit cycling and conditions between DPWM and ADC resolution for its elimination \cite{peterchev2003,peng2007}, while quantization-effect modeling permits quantitative prediction of these oscillations \cite{maksimovic2007}. Recent work extends the analysis to MHz-range GaN converters using dither and sigma-delta modulation \cite{alobaidi2025,gonzalez2020} and to design guidelines for two-loop current-mode digital controllers \cite{gonzalez2020}.

Digital control relieves the analog constraints of passive compensation, enables programmability and auto-tuning, and integrates naturally with digital loads \cite{maksimovic2004,corradini2015}. In exchange, the designer accepts three implementation-quantized quantities that an analog loop never sees: the ADC resolution of the measured output, the DPWM resolution of the commanded duty, and the sampling period at which controller and power stage are synchronized.

The theoretical basis is well established. Peterchev and Sanders derived conditions on ADC/DPWM resolution trade-offs for limit-cycle elimination and introduced digital dither to recover effective DPWM resolution \cite{peterchev2003}. Peng, Prodić, Alarcón, and Maksimović derived static and dynamic quantization models explaining the origins of limit-cycle oscillations \cite{peng2007}, and Maksimović and Zane provided the small-signal discrete-time framework used to design digitally controlled converters \cite{maksimovic2007}. The broader landscape is surveyed in \cite{liu2009}. On the application side, low-resolution non-uniform ADC quantization has been shown to improve dynamic response \cite{hu2008}, and the full control loop benefits from predictive structures \cite{chen2003,bibian2002}.

Two observations motivate this work. First, most quantization studies isolate a single stage, whereas a practical controller is constrained in all three quantities simultaneously; the interaction between sampling period, ADC word length, and PWM word length is rarely treated as one design problem. Second, the buck-boost converter, whose inverted output and non-minimum phase control path \cite{yousefzadeh2007} make regulation inherently harder, is underrepresented in the digital-control literature: most buck-boost-specific digital studies consider time-domain or PID structures at a single resolution \cite{vijayalakshmi2014,viswanatha2020} and do not sweep the full implementation space. The classical converter model itself is textbook material \cite{erickson2001}.

This paper reports a controlled simulation study of a digitally controlled inverting buck-boost converter under simultaneous quantization of sampling period, ADC resolution, and PWM resolution. Contributions: (1) a 79-run factorial study that separates the individual effect of each quantizer from their combined effect and locates the critical sampling boundary; (2) an objective worst-case selection made after the sweep, not a priori; (3) the identification of a half-scale ADC coincidence ($V_{ref} = -FS/2$) that masks steady-state error at low ADC resolutions; (4) formal evidence that interactions between the three stages amplify beyond additivity, quantified by an interaction index; (5) design guidance linking sampling period and resolution choices, summarized as a performance map; and (6) operating-point robustness checks confirming the conclusions away from the nominal point. Section II develops the converter and controller models; Section III defines the experimental protocol; Section IV presents results; Section V discusses design implications; Section VI concludes.


\section{System Model and Digital Control}

The plant is an ideal inverting buck-boost converter with input $V_{in}=+12$\,V, reference $V_{ref}=-12$\,V, $L=100\,\mu$H, $C=220\,\mu$F, switching frequency $f_{sw}=100$\,kHz ($T_{sw}=10\,\mu$s), and load $R=10\,\Omega$, modeled after the inverting topology whose steady-state point in continuous conduction mode (CCM) is $|V_C|=D/(1-D)\,V_{in}=12$\,V at $D=0.5$, $i_L=2.4$\,A, $\Delta i_L=600$\,mA. The switched-state model integrates $dt=T_{sw}/200$:

\begin{equation}
\text{ON:}\quad \frac{di_L}{dt}=\frac{+V_{in}}{L}, \quad \frac{dv_C}{dt}=-\frac{v_C}{RC}
\end{equation}

\begin{equation}
\text{OFF:}\quad \frac{di_L}{dt}=\frac{v_C}{L}, \quad \frac{dv_C}{dt}=-\frac{i_L+v_C/R}{C}
\end{equation}

During OFF, the inductor current is held at zero when it would otherwise reverse (zero-current boundary), which counts discontinuous-mode (DCM) cycles and keeps the model valid in DCM without clamping the output---a point on which many simplified studies go silent. All reported results are obtained on this switched-state model, not on an averaged approximation, so quantization artifacts (limit cycles, DCM chatter) appear faithfully.

\subsection{Digital Controller and Quantization}

A velocity-form digital PID with positive gains,

\begin{equation}
d[k] = d[k-1] + K_p\big(e[k]-e[k-1]\big) + K_i T_s\, e[k] + \frac{K_d}{T_s}\big(e[k]-2e[k-1]+e[k-2]\big)
\end{equation}

with $e[k]=v_{meas}[k]-V_{ref}$, duty saturated to $[0.05,\,0.95]$, and no separate integrator state (inherent anti-windup). The measurement, duty, and update cadence are each quantized independently:

\begin{itemize}
\item Sampling: the controller updates every $T_s \in \{T_{sw},\,2T_{sw},\,5T_{sw}\}$ (10, 20, 50 $\mu$s);\\ ADC: mid-tread quantizer, full-scale $FS=24$\,V, 12/10/8/6 bits. Note $V_{ref}=-FS/2$, so the reference is exactly representable at every resolution;\\ PWM: duty quantized to $2^b$ levels, $b\in\{12,10,8,6\}$.
\end{itemize}

A continuous-ideal baseline recomputes the same PID every $dt$ with no quantization (true analog reference). The frozen controller gains were obtained once, deterministically, before any quantization study: grid search over $K_p\in\{0.005,0.01,0.02,0.05\}$, $K_i\in\{20,100,500,2000\}$, $K_d\in\{0,2\times10^{-6},10^{-5}\}$ with cost $J=o_v/10+t_s/5\,\text{ms}+\text{IAE}/0.05$ under constraints overshoot $\le 10\%$ and $|e_{ss}|\le 60$\,mV, re-validated on the switched model: $K_p=0.05$, $K_i=100.0$, $K_d=10^{-5}$. No parameter was touched during the quantization study.



\section{Simulation Framework}

The study executes 79 simulations (Table~\ref{tab:plan}): the continuous baseline, the ideal digital controller ($T_s=T_{sw}$, 16-bit ADC, 16-bit PWM), three sampling-only points, a six-point focused sampling sweep across the critical boundary, four ADC-only, four PWM-only runs, the full $3\times4\times4$ factorial sweep (48 runs), six disturbance runs (input step $12\to9$\,V and load step $10\to5\,\Omega$, applied at $t=60$\,ms, on the continuous baseline, the ideal digital controller, and the objectively selected worst configuration), and six operating-point robustness runs ($V_{in}=10,15$\,V on three configurations).

\begin{table*}[htbp]
\centering
\begin{tabular}{l l l l}
\toprule
Group & Runs & T\_s & ADC/PWM \\
\midrule
Continuous baseline & 1 & --- & analog \\
Ideal digital & 1 & 1x & 16/16 bit \\
Sampling only & 3 & 1/2/5x & 16/16 \\
Boundary sweep & 6 & 1.5-4x & 16/16 \\
ADC only & 4 & 1x & 12/10/8/6, 16 \\
PWM only & 4 & 1x & 16, 12/10/8/6 \\
Factorial sweep & 48 & all & all \\
Disturbances & 6 & selected & 3 configs \\
Op-point robustness & 6 & selected & 3 configs \\
\bottomrule
\end{tabular}
\caption{Study layout (79 simulations).}
\label{tab:plan}
\end{table*}

Metrics: 2\textbackslash{}\% settling time, overshoot, steady-state error, peak-to-peak ripple (limit cycle), IAE, ITAE, and DCM-cycle count. For disturbance runs all metrics are evaluated on the post-step 60\,ms window and ripple on its final 50\textbackslash{}\%. The worst configuration is selected objectively after the sweep, as the run with the largest normalized degradation of steady-state error, IAE, and limit cycle relative to the ideal digital controller---no configuration is pre-declared worst.


\section{Results}

\subsection{Continuous versus Ideal Digital}

The conversion from analog to ideal digital control already costs measurable quality (Table~\ref{tab:base}): the 16-bit/one-cycle system settles 36\textbackslash{}\% slower (2.66 vs 1.96\,ms) and shows a 12.76\,mV steady-state offset and 39.6\,mV limit cycle against the continuous 0.44\,mV / 26.7\,mV. The 12.8\,mV floor is the quantization-limited operating point for these gains; every further result is expressed relative to it.

\begin{table*}[htbp]
\centering
\begin{tabular}{l l l}
\toprule
 & Continuous & Ideal digital \\
\midrule
Settling (ms) & 1.96 & 2.66 \\
Overshoot (\%) & 0.07 & 0.02 \\
|e\_ss| (mV) & 0.44 & 12.76 \\
Limit cycle (mV) & 26.7 & 39.6 \\
IAE & 0.0058 & 0.0109 \\
DCM cycles & 0 & 0 \\
\bottomrule
\end{tabular}
\caption{Continuous baseline vs ideal digital.}
\label{tab:base}
\end{table*}


\subsection{Individual Stage Effects}

Fig.~\ref{fig:sampling} isolates the sampling stage at 16-bit resolution and locates the critical boundary. Settling degrades smoothly as $T_s/T_{sw}$ rises from 1 to 3.5 (2.66 to 6.38 ms, all still regulation-valid), then the loop loses regulation outright at $4\times$: it never settles, the limit cycle jumps to 1.59 V, and 101\,736 DCM cycles accrue in 60 ms. The practical boundary therefore sits between $3.5\times$ and $4\times T_{sw}$ (35--40 $\mu$s); $T_s = 5T_{sw}$ fails in every configuration, alone or combined.

\begin{figure*}[htbp]
\centering
\includegraphics[width=0.8\textwidth]{/home/touhid/Documents/powerpaper/results/figures/fig2_sampling.png}
\caption{Sampling-period boundary sweep at 16-bit ADC/PWM: settling time, steady-state error, limit cycle, and IAE versus $T_s/T_{sw}$. Regulation is lost between $3.5\times$ and $4\times T_{sw}$.}
\label{fig:sampling}
\end{figure*}

ADC resolution couples to the reference in a subtle way (Fig.~\ref{fig:adc}): because $V_{ref}=-FS/2$ is exactly representable, the 8-bit and 6-bit ADCs lock onto the exact code and reduce the measured steady-state error to 0.44\,mV while masking a worse transient (overshoot rises from 0.06\textbackslash{}\% at 12-bit to 0.85\textbackslash{}\% at 6-bit; settling from 2.70 to 3.49\,ms). This is a half-scale coincidence, not an ADC benefit.

\begin{figure*}[htbp]
\centering
\includegraphics[width=0.8\textwidth]{/home/touhid/Documents/powerpaper/results/figures/fig3_adc.png}
\label{fig:adc}
\end{figure*}

PWM resolution below 8 bits dominates the limit cycle (Fig.~\ref{fig:pwm}): 6-bit PWM adds a 204.8\,mV ripple versus $\approx$39.5\,mV at 12/10/8-bit. Duty quantization is the direct source of steady-state ripple, as predicted by limit-cycle analysis \cite{peterchev2003,peng2007}; at 6-bit the DPWM step is $2/64=31.25$\,mV of duty, which maps through the control-to-output gain into the observed 204.8\,mV ripple. This is precisely the regime addressed by finer DPWM controller implementations \cite{patella2003,syed2004} and by noise-shaping modulators \cite{norris2008}.

\begin{figure*}[htbp]
\centering
\includegraphics[width=0.8\textwidth]{/home/touhid/Documents/powerpaper/results/figures/fig4_pwm.png}
\label{fig:pwm}
\end{figure*}


\subsection{Combined Sweep and Interactions}

In the 48-run factorial sweep the worst configuration is $T_s=5T_{sw}$, ADC=8-bit, PWM=12-bit (normalized degradation score 1309.2): settling never completes within 60\,ms, steady-state error $-13.43$\,V, limit-cycle amplitude 7.33\,V, IAE 0.789, 127\,075 DCM cycles. Fig.~\ref{fig:perf} classifies every combination against explicit thresholds (settling $\le 6$\,ms, $|e_{ss}| \le 120$\,mV, limit cycle $\le 120$\,mV, overshoot, DCM) into excellent / acceptable / marginal / unusable. Sampled at $T_s = T_{sw}$, 13 of 16 combinations are excellent and the rest acceptable; at $2\times$ all are acceptable or marginal; at $5\times$ all 16 are unusable.

\begin{figure*}[htbp]
\centering
\includegraphics[width=0.8\textwidth]{/home/touhid/Documents/powerpaper/results/figures/fig8_perf_map.png}
\caption{Performance map of the 48-run sweep. Thresholds: excellent ($t_{settle}\le4$ ms, $|e_{ss}|\le60$ mV, LC $\le60$ mV, ov $\le5\%$), acceptable ($t_{settle}\le6$ ms, $|e_{ss}|,\text{LC}\le120$ mV), marginal, unusable ($t_{settle}\ge50$ ms).}
\label{fig:perf}
\end{figure*}

The worst configuration is far worse than any individual degradation: $-13.43$\,V versus the $-13.07$\,V of sampling alone, but with 127\,075 DCM cycles and a 7.33\,V limit cycle that the sampling-only run does not exhibit. The super-additivity is quantitative: with $D$ the normalized partial degradation (steady-state error, IAE, and limit cycle relative to the ideal digital controller), the individual terms are $D_{sampling}(5\times)=1141.6$, $D_{ADC}(8\text{-bit})=0.6$, and $D_{PWM}(12\text{-bit})=0.0$; the additive prediction is 1142.2, whereas the measured combined degradation is $D_{total}=1309.2$---an interaction of $+167$ ($+15\%$). The 8-bit ADC coincidence, harmless alone, converts a slowly sampling loop into deep limit-cycle chatter once the 12-bit PWM pulls the operating point across the coincidence.


\subsection{Disturbance Rejection}

Table~\ref{tab:dist} summarizes the six disturbance runs (step at $t=60$\,ms; recovery time $=t_{settle}-60$\,ms). Both healthy controllers reject input and load steps in under 1.5\,ms with negligible IAE. The worst configuration never re-regulates within the observation window: it settles into deep DCM chatter with $>$5\,V of steady-state error and over 160\,000 DCM cycles per event (Fig.~\ref{fig:dist}).

\begin{table*}[htbp]
\centering
\begin{tabular}{l l l l}
\toprule
Config & Event & Recovery & |e\_ss| / IAE \\
\midrule
Continuous & Vin step & 1.43 ms & 0.0 mV / 0.0014 \\
Continuous & Load step & 0.94 ms & 0.3 mV / 0.0013 \\
Ideal dig. & Vin step & 0.92 ms & 14.9 mV / 0.0017 \\
Ideal dig. & Load step & 0.84 ms & 26.1 mV / 0.0019 \\
Worst & Vin step & never & 8.9 V / 0.535 \\
Worst & Load step & never & 5.9 V / 0.353 \\
\bottomrule
\end{tabular}
\caption{Disturbance response (step at $t=60$\,ms).}
\label{tab:dist}
\end{table*}

\begin{figure*}[htbp]
\centering
\includegraphics[width=0.8\textwidth]{/home/touhid/Documents/powerpaper/results/figures/fig7_disturbances.png}
\label{fig:dist}
\end{figure*}


\subsection{Operating-Point Robustness}

The main conclusions are not an artifact of the single nominal point ($V_{in}=12$\,V, $V_{ref}=-12$\,V, $R=10\,\Omega$). Table~\ref{tab:op} repeats the ideal, a representative good ($1\times$, 10-bit, 10-bit), and a representative poor ($2\times$, 6-bit, 6-bit) configuration at $V_{in}=10$ and 15\,V with the same frozen gains. Ideal and good configurations retain full regulation (2.5--2.7\,ms settling, sub-100\,mV limit cycle) across the input range. The poor configuration holds at $10$ and $12$\,V but degrades sharply at $V_{in}=15$\,V---it settles only at the very edge of the 60\,ms window (59.3 ms) with a 561\,mV limit cycle and 18 DCM excursions---showing that marginal designs are also operating-point sensitive, not just sampling-sensitive.

\begin{table*}[htbp]
\centering
\begin{tabular}{l l l l}
\toprule
config & 10 V & 12 V & 15 V \\
\midrule
ideal (1×,16,16) & 2.74 / 67.9 & 2.66 / 39.6 & 2.55 / 28.9 \\
good (1×,10,10) & 2.69 / 28.6 & 2.64 / 57.8 & 2.52 / 80.6 \\
poor (2×,6,6) & 5.40 / 28.9 & 5.11 / 27.7 & 59.3 / 561* \\
\bottomrule
\end{tabular}
\caption{Operating-point robustness: settling time (ms) / limit cycle (mV) at $V_{in}=10,12,15$\,V with gains frozen ($V_{ref}=-12$\,V, $R=10\,\Omega$). No asterisk: settles inside 60 ms; * does not settle in the window.}
\label{tab:op}
\end{table*}



\section{Discussion}

The results position the three quantizers on a clear hierarchy for this converter: sampling period dominates (a 50\,$\mu$s loop is unstable regardless of resolution); DPWM resolution sets the achievable steady-state ripple; ADC resolution interacts with the reference placement. This ordering matches the known limit-cycle conditions: DPWM quantization bounds the smallest controllable output step, and the controller can only hold the output within half an ADC code of the reference \cite{peterchev2003}. Our 12.8\,mV ideal-digital floor is close to the $FS/2^{b+1} = 24/8192 \approx 2.9$\,mV ADC quantum times the loop gain, consistent with the static quantization model of \cite{peng2007}.

The half-scale coincidence ($V_{ref} = -FS/2$) warrants a design caveat: converters with mid-rail references silently benefit from exact code lock-on and misreport their true regulation quality. A $V_{ref}$ placed off half-scale, or a smaller ADC full-scale, removes the coincidence and exposes the 8-bit/6-bit error directly. Non-uniform ADC quantization \cite{hu2008} offers one systematic remedy by assigning finer codes near the reference instead of relying on coincidence.

Where the literature targets specific remedies, our sweep quantifies the room they have to work: sigma-delta and dither techniques \cite{alobaidi2025,lukic2007} recover effective PWM resolution but cannot rescue an unstable sampling rate; predictor-based controllers \cite{chen2003,bibian2002} and auto-tuning of limit cycles \cite{zhao2007} act on the ADC side and are complementary to the sampling constraint identified here; DPWM delta-sigma architectures \cite{norris2008,syed2004} similarly address only the ripple term. The engineering takeaway is that sampling-period choices are a prerequisite, and resolution techniques are refinements on top. As a practical rule from the performance map, sampling at or below $3.5\times T_{sw}$ with mid-range resolutions keeps regulation inside the acceptable band; the combination of a slow sample rate with a coarse quantizer is the only regime that is unusable in the entire map \cite{peng2007}.


\section{Conclusion}

Sampling period is the dominant quantization-stage effect for this converter: a 200\,$\mu$s loop period is unstable alone and in every combination, and no amount of ADC or PWM resolution rescues it. The critical boundary is sharper than the discrete study implied: regulation degrades smoothly to $3.5\times T_{sw}$ and fails between $3.5\times$ and $4\times$. ADC and PWM resolution contribute distinct, additive-but-interacting artifacts, and the paper's key design guidance is that resolution choices cannot be made independently of the sampling period. The ideal digital controller captures 95\textbackslash{}\% of continuous-loop performance; the remaining gap is the quantization floor. The main conclusions are robust to the operating point, although marginal (coarse) configurations show additional operating-point sensitivity. All figures and 79-run data are available as simulation artifacts alongside this paper.

Future work will extend the sweep to CCM/DCM boundary load points and to cascaded current-mode structures \cite{gonzalez2020}, where the inner current loop is itself subject to sampling and quantization effects, and will validate the half-scale coincidence hypothesis experimentally on a prototype.


\section{Acknowledgment}

The authors thank Tareq Anwar Shikdar for guidance; no external funding was received.


\bibliographystyle{IEEEtran}
\bibliography{references}

\end{document}